\chapter{Running Akida models}
\label{chap:operacion}

\section{Load and inspect a model}

Akida models are distributed as FBZ files. Before running one, check at least its
input shape, output shape and the preprocessing used during conversion.

Load the model and inspect its structure:

\begin{lstlisting}[language=Python]
import akida

model = akida.Model("models/model.fbz")

print("input_shape:", model.input_shape)
print("output_shape:", model.output_shape)
print("ip_version:", model.ip_version)
model.summary()
\end{lstlisting}

\code{summary()} shows the layers and how the model can run. Resource use depends
on the model and mapping mode.

\section{Select the device}

\code{akida.devices()} returns the available physical devices \cite{brainchip_akida_guide}:

\begin{lstlisting}[language=Python]
import akida

devices = list(akida.devices())

if not devices:
    raise RuntimeError("No Akida device detected")

device = devices[0]
print(device)
\end{lstlisting}

A device created with \code{akida.AKD1000()} is virtual and checks model mapping.
To run inference on the board, use a device returned by \code{akida.devices()}.

\section{Map the model to the board}

A loaded model can initially run in software. To use the AKD1000, map it to the device:

\begin{lstlisting}[language=Python]
model.map(
    device,
    hw_only=True,
    mode=akida.MapMode.AllNps
)

model.summary()
\end{lstlisting}

\code{hw\_only=True} requires execution entirely on Akida hardware
\cite{brainchip_akida_guide}. If the model cannot map this way, the call raises an error.

\code{AllNps}, \code{HwPr} and \code{Minimal} allocate AKD1000 resources in different
ways. For a first run, you can use \code{AllNps}. \Cref{chap:rendimiento} examines
its effect on latency and power consumption.

\section{Prepare the input}

Provide the input as a NumPy \code{uint8} tensor. Its shape, effective range and
channel order must match those used during model conversion. Check each sample
shape against \code{model.input\_shape}.

For example, for data stored in NumPy:

\begin{lstlisting}[language=Python]
import numpy as np

data = np.load("data/inputs.npy", allow_pickle=False)

print("shape:", data.shape)
print("dtype:", data.dtype)
print("min/max:", data.min(), data.max())

if data.dtype != np.uint8:
    raise TypeError("Akida model inputs must use uint8")

if tuple(data.shape[1:]) != tuple(model.input_shape):
    raise ValueError(
        f"Input shape {data.shape[1:]} "
        f"!= {model.input_shape}"
    )
\end{lstlisting}

The effective range depends on the model input precision. The tool rejects other
data types without automatically converting them to \code{uint8}.


\section{Run the model}

The Akida API can run the model through \code{forward()} or \code{predict()}.
The choice depends on how the model output was defined and validated during conversion
\cite{brainchip_akida_guide}.

For example:

\begin{lstlisting}[language=Python]
output = model.forward(data)
print(output)
\end{lstlisting}

Keep the same execution method throughout a measurement campaign so that the results
are comparable.

\pagebreak
You can also use the tool included in this repository:

\begin{lstlisting}[language=bash]
akd1000-bringup run \
  --model models/model.fbz \
  --input data/input.npy \
  --method forward \
  --clock Performance \
  --map AllNps \
  --output outputs/raw-output.npy
\end{lstlisting}

The tool loads the input, runs the model and saves the output to a NumPy file.
The input must have the shape

\begin{lstlisting}
(batch, *model.input_shape)
\end{lstlisting}

Models with multiple inputs or outputs require direct handling through the Akida API.

\section{Clock mode}

The AKD1000 has \code{Performance}, \code{Economy} and \code{LowPower} clock modes.
Select the mode through the API:

\begin{lstlisting}[language=Python]
device.soc.clock_mode = akida.soc.ClockMode.Performance
print(device.soc.clock_mode)
\end{lstlisting}

Record the clock mode when comparing performance measurements. Lower instantaneous
power does not necessarily mean lower energy per inference, because execution time
may also change.
